\section{Discussion}
\label{sec:discussion}

What happens in a failed run is easy to misread. Loss falls smoothly, so it
is not divergence; the val--test gap is small, so it is not overfitting; and
it is not the attention problem reported for vision Transformers
\cite{zhai2023stabilizing}, where attention sharpens too early. Here it
flattens until every time step is weighted equally. The network has found
the simplest way to reduce MSE, which is to output the label mean, and then
never leaves before the learning-rate budget and early-stopping patience run
out. The failure is hard to spot because that answer is moderately good
under the training loss; only the score used in service and the checks of
Section~\ref{sec:method_audit} reveal it.

\subsection{Limitations}

All quantitative statements come from C-MAPSS. The explanation is framed in
terms of architecture and data structure, but whether it carries over to
N-CMAPSS, milling or bearing run-outs is a hypothesis, not a result. Rules
1--2 were read off one set of experiments rather than derived from theory,
and the FD004 result shows that local rules cannot guarantee reliable
training even within one benchmark. Five seeds per configuration is more
than is usually reported, but it is still a coarse estimate of a failure
probability near 0 or 1; the 6/8 FD004 rate has a wide binomial interval.
Our audit holds the train/val split fixed across seeds, matching the
archived runs, while the instrumented study also varied it and outcomes
changed (FD004, seed 123: failed under the fixed split, worked under the
varied split); the ``seed'' whose reliability one audits should include the
split. Windowed train/val splits from the same trajectories also inflate
validation optimism relative to engine-holdout splits, at the cost of
val--test gaps up to 25\% on FD003.

\subsection{Practice and outlook}

The checks are not tied to our architecture. Report mean $\pm$ std over at
least three seeds together with the failure rate, and treat best-of-$k$
tuning as suspect by default. Record two numbers during training, prediction
std on a held-out batch and mean attention entropy, whose thresholds are
sharp (0.01 cycles; 0.99); watch for a validation RMSE that sits at
$\operatorname{std}(y_{\mathrm{val}})$ with near-constant predictions, which
means the run is still printing the average and should be started again with
a new seed or given more patience. Keep more than one usable temporal branch,
since removing one for the sake of a smaller model should require evidence
that training is unaffected, not just that accuracy is; and judge the
deployed model on the asymmetric score. If recovery depends on a jump in the
TCN gradients, warm-up schedules or a learning-rate increase triggered by a
long flat stretch might make the seed-dependent settings train reliably; the
per-epoch records give a way to test that, and the same checks apply to any
sensor-regression task with several operating conditions.

\section{Conclusion}
\label{sec:conclusion}

We documented, and partly fixed, a training failure in the hybrid deep
models used for aero-engine RUL prediction on multi-condition data. Across
161 runs, 37.2\% of FD002 and 24.2\% of FD004 runs ended by predicting one
fixed number for every engine, whose test error can be computed in closed
form; none of 81 single-condition runs did. Every run first spends time
printing the training-set average, with features organised by operating
condition and attention weights identical at every time step. Training
starts to improve again only when gradients in the convolutional branch jump
about a hundredfold, and this has to happen before early stopping. The rules
we propose choose window length and depth by the number of operating
conditions, keep several usable branches, drop the adversarial and
contrastive components, and add three cheap checks; over five seeds they
trained reliably on FD001/FD002/FD003 (13.66/19.43/16.09 RMSE), while on
FD004 six of eight runs still failed and the two that worked agreed at
19.53 $\pm$ 0.07, which the checks catch in time to start again. A table of
single-run accuracies cannot say whether a model is fit for safety-critical
use: one that may, on retraining, quietly turn into a device printing the
fleet-average RUL for every engine is not usable at any RMSE. That
probability has almost never been measured; we supply the measurement, the
explanation and the rules.
